% ============================================================
% Chapter 7: Results
%   Section: The Stages Before the Verdict
% Included from chapters/07-results/results.tex
% Source: docs/experiments.md; backend/data/evaluation/experiments/;
% docs/decisions/performance.md (time per stage, the 20-second contest,
% 2026-10-20 to 22); experiments/speed/rule.md.
% Discovery and the candidate ranking are reported under RQ1
% (Section~\ref{sec:results-rq1}); this section holds the rest.
% ============================================================

\section{The Stages Before the Verdict}
\label{sec:results-pre-verdict}

Discovery and the candidate ranking, measured on 5, 6, 9 and
13~October, answer RQ1 and are reported in
Section~\ref{sec:results-rq1}. This section reports the other stages
that decide which articles get a verdict.

\paragraph{The mission screen.} On 527 live candidates from Culture,
Society and Health rounds (6~October), each title and summary was
compared with the topics and the off-mission descriptions exactly as
the screen does. At a margin of 0.04 it left out 68 of the 484
candidates from general outlets; read by hand, nearly all were
elections, crime, accidents, storms, market wire items and celebrity
interviews, the doubtful ones four war-history pieces and one story on
data-centre regulation. Between margins of 0.03 and 0.04 about 3 in 17
candidates were worth keeping, which is why the line is at 0.04.
Widened the same day with descriptions for labour disputes and
markets, it left out 115 of 545 instead of 80, all 36 new ones being
strikes, protests, finance and election politics, with one story wrongly
let back in. The sentiment model was tried as a signal and rejected: it
rated an endangered primate's birth 0.89 negative and a lost theatre
reborn 0.89 negative, because a problem being solved reads as negative.

\paragraph{Admission.} Replaying the admission thresholds offline over
the articles of the 41 facts labelled by 5~October, 71\% would be
admitted at the values of the local configuration that day (topic
minimum 0.35, impact minimum 0.3, against a declared default of 0.1 for
the latter); at a topic minimum of 0.40, 61\% at an impact minimum of
0.1 and 49\% at 0.3. Four articles (10\%) have no topic at the
classifier's floor and are rejected whatever the minimum.
\pending{the admission experiment over at least 100 round candidates
the editor chose for or against: admission rate, precision and recall
over the grid, and the AUC of the positive-impact score.}

\paragraph{The topic classifier.} On the same 41 articles, the
classifier's top topic was the gold one 54\% of the time [37, 68], the
gold topic was in its top three 78\% [66, 90], its top group was right
73\% [59, 85], and 10\% [2, 20] of articles had no topic at the 0.35
floor. Society is the weak group: five of its nine articles were
classified as environment. Topics steer admission and the balance of
the set, not verdicts, which is why improving the classifier was
de-scoped on 29~September.

Run again on 20~October over the 58 articles of the custom set's facts,
with the keyword lists widened since: the top topic was the gold one
67\% of the time [55, 79], the gold topic in its top three 88\% [79, 95],
its top group right 83\% [72, 91], and 2\% [0, 5] of articles had no
topic at the floor --- so the 0.35 floor stays. Society is still the weak
group: five of its nine articles right, the other four spread over the
other groups; three of seven culture articles went to environment.
\pending{the classifier on the reviewed labels: the topics are the AI
pass's for most of these articles.}

\paragraph{The sources, week by week.} A weekly check runs the source
probe and joins it with the time to reception. Its first week,
20~October: 30 of 33 sources up, 3 degraded and none down; 157 of 165
sample articles extracted (95.2\%), every one by the first step of the
cascade, so the later steps saved nothing that day. El~País refused
three of five politics pages with HTTP 403, and the headless browser did
not get past it either; National Geographic and SINC have no working
feed and are found through their section pages. Both meta-search checks
answered (86 and 90 results), with Google and Brave suspended by their
own rate limits.
\pending{weeks two to four, and the time to reception per source once
rounds are analysed again.}

\paragraph{Time per stage.} The X-Fact pilot gives the claim-level
split on the development laptop (Section~\ref{sec:results-rq4}); the
production stack's single analysis gives the processor's. One analysed
article of four claims on the development laptop, from the job journal
of 20~October: fetching the article 1~s; enrichment 8--26~s, most of it
the embedding model over the whole text; claim selection 3.5~s; search
and page fetches, every claim at once, 15--30~s; ranking 5--8~s a
claim; the verifier 30--65~s a claim, one call at a time; the graph
write 4~s. Per claim, on the test half on 22~October (the desktop
processor, the list verifier): the claim's entities 2.2~s, retrieval
11.3~s (2.3 of them the web searches, 3.1 the page fetches), ranking
5.5~s and the model 44.9~s at the median, 12.6 of them waiting for the
other claim's call. The model is 70\% of a claim's time, and of one
call most is reading the prompt, not writing the answer: 29~s reading
and 16~s writing for a verifier-shaped call on the laptop.

\paragraph{A twenty-second claim.} Asked to check a claim in about
twenty seconds without a paid model, four configurations were compared
on ten tuning claims by a rule written before the runs, and the same ten
were checked by a hosted assistant (Claude) with web search, blind to
the labels, as a reference. A shorter answer format took the model's
call from 28 to 19~s; fewer and shorter sources, to 8~s; choosing the
pages by the claim's terms and ranking passages rather than page heads
took retrieval and ranking from 16 to 4~s. The fastest profile, on
search snippets alone, took 14.2~s a claim at the median and never more
than 22~s; the assistant took 13.1~s. On the ten claims every
configuration was right 2 or 3 times and the assistant 7. On all 69
tuning claims the same day both fast profiles were measurably less
accurate than the thorough pipeline: 33.3\% and 33.8\% against 50.7\%,
$-17.4$ [$-30.4$, $-4.3$] and $-16.2$ [$-27.9$, $-4.4$] points. The
application stays thorough, with the fast profiles one setting away:
about 15~s a claim at some 17 points of accuracy, or about a minute at
50.7\%. Speed did not cause the 3B model's lean to \texttt{FALSE} and
does not fix it.
\pending{admission's own time, and one ingestion round end to end, from
the metrics endpoint and the job journal.}
